\documentclass[11pt,a4paper]{article}
\input{../template/utm-lab}

\newcommand{\labnumber}{9}
\newcommand{\labtitle}{Watching System Calls with strace}
\newcommand{\studentname}{Daniel Nicorici}
\newcommand{\studentgroup}{FAF-243}

\begin{document}

\utmtitlepage
\tableofcontents
\newpage

\plainsection{Purpose of the work}

A system call is the door between a program and the kernel. The purpose of this lab is to watch
that door with \code{strace}, which prints every system call a program makes: first the calls of
a real command, \code{ls -l}, then the calls of a small program of my own that uses \code{open},
\code{write} and \code{read} directly, and finally to count one kind of call.

The work was done on my Arch Linux laptop with gcc 16.2.1 and strace 7.3, in the folder
\code{\textasciitilde/os-lab9}. In the listings, long runs of similar lines are shortened and
marked.

\section{Part 1. See the calls of a real command}

\begin{lstlisting}[caption={The calls of ls -l (shortened)}]
$ strace -f -e trace=openat,read,write,close ls -l
openat(AT_FDCWD, "/etc/ld.so.cache", O_RDONLY|O_CLOEXEC) = 3
close(3)                                = 0
openat(AT_FDCWD, "/usr/lib/libcap.so.2", O_RDONLY|O_CLOEXEC) = 3
read(3, "\177ELF\2\1\1\0\0\0\0\0\0\0\0\0\3\0>\0\1\0\0\0\320{\0\0\0\0\0\0"..., 1024) = 1024
close(3)                                = 0
openat(AT_FDCWD, "/usr/lib/libc.so.6", O_RDONLY|O_CLOEXEC) = 3
read(3, "\177ELF\2\1\1\3\0\0\0\0\0\0\0\0\3\0>\0\1\0\0\0\300y\2\0\0\0\0\0"..., 1024) = 1024
close(3)                                = 0
[... locale files ...]
openat(AT_FDCWD, ".", O_RDONLY|O_NONBLOCK|O_CLOEXEC|O_DIRECTORY) = 3
openat(AT_FDCWD, "/etc/passwd", O_RDONLY|O_CLOEXEC) = 4
read(4, "root:x:0:0::/root:/usr/bin/bash\n"..., 4096) = 1715
close(4)                                = 0
openat(AT_FDCWD, "/etc/group", O_RDONLY|O_CLOEXEC) = 4
read(4, "root:x:0:root\nsys:x:3:bin\nmem:x:"..., 4096) = 1004
close(4)                                = 0
[... user database lookups ...]
close(3)                                = 0
write(1, "total 4\n", 8)                = 8
openat(AT_FDCWD, "/etc/localtime", O_RDONLY|O_CLOEXEC) = 3
read(3, "TZif2\0\0\0\0\0\0\0\0\0\0\0\0\0\0\0\0\0\0\20\0\0\0\20\0\0\0\0"..., 4096) = 2424
close(3)                                = 0
write(1, "-rw-r--r-- 1 Nico Nico 1055 Nov "..., 52) = 52
close(1)                                = 0
close(2)                                = 0
+++ exited with 0 +++
\end{lstlisting}

\begin{lstlisting}[caption={Counting the calls of one ls -l}]
$ strace -f -e trace=openat,read,write,close ls -l 2>&1 >/dev/null | grep -c '^[a-z]*('
101
$ strace -c ls -l > /dev/null
% time     seconds  usecs/call     calls    errors syscall
------ ----------- ----------- --------- --------- ----------------
 12.71    0.000053           1        41           close
  7.19    0.000030           0        44        13 openat
  4.56    0.000019           0        34           fstat
  3.84    0.000016           1        15           read
  3.36    0.000014           0        36           mmap
  2.16    0.000009           2         6           getdents64
  0.48    0.000002           2         1           write
[... 35 other system calls ...]
------ ----------- ----------- --------- --------- ----------------
100.00    0.000417           1       314        27 total
\end{lstlisting}

The burst at the beginning is not \code{ls} itself yet: the dynamic loader opens the shared
libraries, then the C library opens the locale files. \code{ls} then opens the directory
\code{.}, reads \code{/etc/passwd} and \code{/etc/group} to turn the numeric owner and group
into names, reads \code{/etc/localtime} to print the dates, and only then writes.

\textbf{Observations:}
\begin{enumerate}
  \item \textbf{How \code{ls} prints.} With the \code{write} system call, to descriptor 1,
  standard output: \code{write(1, "total 4\textbackslash n", 8) = 8}, followed by one more
  \code{write(1, ...)} for the line of the file.
  \item \textbf{How many calls.} One simple \code{ls -l} made 101 calls of the four traced kinds
  and 314 system calls in total, 27 of which returned an error (mostly files that do not exist
  in the first place the library looks).
\end{enumerate}

\section{Part 2. A program that calls the kernel directly}

\begin{lstlisting}[style=ccode,caption={syscalls.c}]
#include <fcntl.h>
#include <unistd.h>
#include <string.h>
#include <stdio.h>

int main(void) {
    const char *msg = "hello, kernel\n";
    int fd = open("note.txt", O_CREAT | O_WRONLY | O_TRUNC, 0644);
    if (fd < 0) { perror("open"); return 1; }
    printf("open returned descriptor %d\n", fd);
    ssize_t n = write(fd, msg, strlen(msg));
    printf("write sent %zd bytes\n", n);
    close(fd);
    fd = open("note.txt", O_RDONLY);
    char buf[64];
    n = read(fd, buf, sizeof(buf) - 1);
    buf[n] = '\0';
    printf("read got %zd bytes: %s", n, buf);
    close(fd);
    return 0;
}
\end{lstlisting}

\begin{lstlisting}[caption={Compiling and running syscalls}]
$ gcc -Wall -Wextra -o syscalls syscalls.c
$ ./syscalls
open returned descriptor 3
write sent 14 bytes
read got 14 bytes: hello, kernel
\end{lstlisting}

\begin{lstlisting}[caption={The same program seen with strace}]
$ strace -e trace=openat,write,read,close ./syscalls
openat(AT_FDCWD, "/etc/ld.so.cache", O_RDONLY|O_CLOEXEC) = 3
close(3)                                = 0
openat(AT_FDCWD, "/usr/lib/libc.so.6", O_RDONLY|O_CLOEXEC) = 3
read(3, "\177ELF\2\1\1\3\0\0\0\0\0\0\0\0\3\0>\0\1\0\0\0\300y\2\0\0\0\0\0"..., 1024) = 1024
close(3)                                = 0
openat(AT_FDCWD, "note.txt", O_WRONLY|O_CREAT|O_TRUNC, 0644) = 3
write(1, "open returned descriptor 3\n", 27) = 27
write(3, "hello, kernel\n", 14)         = 14
write(1, "write sent 14 bytes\n", 20)   = 20
close(3)                                = 0
openat(AT_FDCWD, "note.txt", O_RDONLY)  = 3
read(3, "hello, kernel\n", 63)          = 14
write(1, "read got 14 bytes: hello, kernel"..., 33) = 33
close(3)                                = 0
+++ exited with 0 +++
\end{lstlisting}

The first five lines are the setup: the loader opens the C library. After them come my calls.
The C function \code{open} appears as the system call \code{openat}, and each \code{printf}
line appears as a \code{write} to descriptor 1; the program's own text printed between the
strace lines is left out of the figure.

\textbf{Observations:}
\begin{enumerate}
  \item \textbf{Why the descriptor is 3.} Descriptors 0, 1 and 2 are already open when the
  program starts: standard input, standard output and standard error, inherited from the shell
  that started it. The kernel always gives the smallest free number, so the first file the
  program opens gets 3. After \code{close(3)} the number is free again, which is why the second
  \code{openat} also returns 3.
\end{enumerate}

\section{Part 3. Count one kind of call}

The handout I received ends after Part 2, so for this part I counted one kind of call,
\code{write}, in my own program, using the summary mode of strace (\code{-c}).

\begin{lstlisting}[caption={All the calls of syscalls, counted}]
$ strace -c ./syscalls > /dev/null
% time     seconds  usecs/call     calls    errors syscall
------ ----------- ----------- --------- --------- ----------------
  0.00    0.000000           0         2           read
  0.00    0.000000           0         2           write
  0.00    0.000000           0         4           close
  0.00    0.000000           0         8           mmap
  0.00    0.000000           0         4           openat
[... 13 other system calls ...]
------ ----------- ----------- --------- --------- ----------------
100.00    0.000000           0        41         2 total
\end{lstlisting}

\begin{lstlisting}[caption={Counting only write, on the terminal and redirected}]
$ strace -c -e trace=write ./syscalls
open returned descriptor 3
write sent 14 bytes
read got 14 bytes: hello, kernel
% time     seconds  usecs/call     calls    errors syscall
------ ----------- ----------- --------- --------- ----------------
  0.00    0.000000           0         4           write
------ ----------- ----------- --------- --------- ----------------
100.00    0.000000           0         4           total
$ strace -c -e trace=write ./syscalls > /dev/null
% time     seconds  usecs/call     calls    errors syscall
------ ----------- ----------- --------- --------- ----------------
100.00    0.000024          12         2           write
------ ----------- ----------- --------- --------- ----------------
100.00    0.000024          12         2           total
\end{lstlisting}

\textbf{Observations:}
\begin{enumerate}
  \item \textbf{The count.} The whole program makes 41 system calls, although the source has
  only 7; the rest is the start-up of the process. On the terminal there are 4 \code{write}
  calls: one to the file and one for each of the three \code{printf} lines.
  \item \textbf{The count depends on where the output goes.} With the output redirected, there
  are only 2 \code{write} calls. \code{printf} is a library function with a buffer: on a
  terminal it flushes at every newline, but into a file or \code{/dev/null} it keeps the three
  lines and hands them to the kernel in a single \code{write} at exit. My direct \code{write}
  to \code{note.txt} is always exactly one system call.
\end{enumerate}

\plainsection{Conclusions}

With \code{strace} the border between a program and the kernel became visible. Even a simple
\code{ls -l} crosses it 314 times: to load its libraries, to read the directory, to look up
user names and the time zone, and finally to print with \code{write} on descriptor 1.
Everything a program does outside its own memory is a system call.

My own program showed the same from the inside. The C functions \code{open}, \code{write},
\code{read} and \code{close} map directly to the calls \code{openat}, \code{write}, \code{read}
and \code{close}, with the arguments I passed and the values I got back, and the first file
received descriptor 3 because 0, 1 and 2 were already taken. Counting the calls also showed the
difference between a system call and a library function: \code{write} always enters the kernel
once, while \code{printf} decides by itself when to do it.

\end{document}
